% =====================================================================
% ANEXO - RESULTADOS SECUNDARIOS
% =====================================================================
% Por indicacion del director (revision del 6-sep, punto 9): recoger aqui los
% resultados de detalle que sostienen los capitulos sin entorpecer su lectura.
%
% Las tablas NO se escriben a mano: las genera scripts/figuras/gen_anexo.py
% desde resultados.py, igual que el resto de las cifras de la memoria.
% =====================================================================

Este anexo recoge mediciones que sostienen los capítulos~\ref{ch:regimenes}
y~\ref{ch:generalizacion} pero que no hacen falta para seguir su argumento. Se
separan aquí por eso, y no porque sean menos fiables: proceden de las mismas
ejecuciones y del mismo módulo de resultados que las del cuerpo.

\section{El \emph{recall} detrás del macro-F1}

La sección~\ref{sec:marco} declara que se reportan dos métricas, el macro-F1 y
el \emph{recall} de la clase de ataque, y explica por qué la segunda hace falta:
un macro-F1 mediocre puede corresponder tanto a un clasificador que se equivoca
como a uno que no detecta nada, y solo el \emph{recall} distingue los dos casos.
El cuerpo de la memoria recoge el macro-F1 en la tabla~\ref{tab:saturacion}; la
tabla~\ref{tab:anexo-recall} completa el par.

La lectura confirma la saturación por la segunda vía. Ninguna vista deja escapar
más del 3\,\% de los ataques en ningún día, de modo que el diagnóstico del
capítulo~\ref{ch:regimenes} no depende de haber elegido una métrica concreta.
% @fuente [R.SATURACION_RECALL[d][v][0] for d in DIAS for v in VISTAS]
\begin{table}[htbp]
  \centering
  \caption[Recall intra-día por vista]{\emph{Recall} de la clase de ataque en la evaluación intra-día. Complementa la tabla~\ref{tab:saturacion}, que recoge el macro-F1 de las mismas mediciones.}
  \label{tab:anexo-recall}
  \begin{tabular}{lccc}
    \toprule
    \textbf{Día} & \textbf{Payload} & \textbf{Metadatos} & \textbf{Conducta} \\
    \midrule
    14-02 cifrado & 0.9997 & 0.9997 & 0.9997 \\
    22-02 en claro & 0.9867 & 0.9887 & 0.9714 \\
    16-02 volumétrico & 1.0000 & 0.9981 & 0.9998 \\
    15-02 volumétrico & 0.9997 & 0.9876 & 0.9963 \\
    20-02 vol. diluido & 1.0000 & 1.0000 & 0.9984 \\
    02-03 periodicidad & 0.9987 & 0.9974 & 0.9988 \\
    \bottomrule
  \end{tabular}
\end{table}


\section{Área bajo la curva con intervalos}

La tabla~\ref{tab:anexo-auc} recoge el AUC-ROC con intervalo de confianza al
95\,\% para los dos regímenes en que se midió con el conjunto benigno ya
corregido. Es la métrica que la sección~\ref{sec:cross-dia} necesita, porque
separa dos preguntas que el macro-F1 confunde: si el modelo ordena bien los
flujos y si el umbral heredado corta donde debe.

El día volumétrico no aparece porque su medición es anterior al reetiquetado, y
las cifras de aquel escenario no son comparables con las de esta memoria por el
motivo que expone la sección~\ref{sec:benigno}. Reportarlas junto a las demás
sería mezclar dos conjuntos de datos distintos bajo la misma etiqueta, que es
uno de los errores que este trabajo documenta haber cometido.
% @fuente [R.RIGOR[r][m][i] for r in R.RIGOR for m in R.RIGOR[r] for i in (0, 1, 2)]
\begin{table}[htbp]
  \centering
  \caption[AUC-ROC con intervalo bootstrap]{Área bajo la curva ROC con intervalo de confianza al 95\,\% de 2.000 remuestreos. Es la medida independiente del umbral que la sección~\ref{sec:cross-dia} emplea para separar el fallo de representación del fallo de umbral.}
  \label{tab:anexo-auc}
  \begin{tabular}{llc}
    \toprule
    \textbf{Régimen} & \textbf{Vista} & \textbf{AUC-ROC (IC 95\,\%)} \\
    \midrule
    Cifrado (14-02, SSH) & payload byte-CNN & 1.0000 [1.0000, 1.0000] \\
     & payload histograma & 0.9958 [0.9902, 1.0000] \\
     & conducta meta + ráfaga & 1.0000 [1.0000, 1.0000] \\
    \midrule
    En claro (22-02, Web) & payload byte-CNN & 1.0000 [1.0000, 1.0000] \\
     & payload histograma & 0.9949 [0.9843, 1.0000] \\
     & conducta meta + ráfaga & 0.9998 [0.9994, 1.0000] \\
    \bottomrule
  \end{tabular}
\end{table}


\section{Detección no supervisada en términos operativos}

La tabla~\ref{tab:anomalia} del cuerpo reporta el área bajo la curva, que resume
la calidad del orden pero no dice cuántos ataques se detectarían en la práctica.
La tabla~\ref{tab:anexo-5fpr} responde a esa pregunta: qué fracción del ataque
se captura si se acepta un 5\,\% de falsos positivos, que es un presupuesto de
alarmas realista para un analista.

La lectura es más dura que la del área bajo la curva y conviene retenerla. En el
régimen cifrado el análisis de contenido no captura ni un solo ataque con ese
presupuesto, pese a un AUC de $0.6946$, mientras la conducta los captura todos.
En el volumétrico ninguno de los tres sirve. Un AUC intermedio puede
corresponder, por tanto, a un detector inservible en el punto de operación que
importa, y esa es una advertencia sobre cómo se leen las cifras de detección de
anomalías en la literatura.
% @fuente [R.ANOMALIA_DET_5FPR[r][m] for r in R.ANOMALIA_DET_5FPR for m in R.ANOMALIA_DET_5FPR[r]]
\begin{table}[htbp]
  \centering
  \caption[Detección no supervisada al 5\,\% de falsos positivos]{Fracción de ataques detectados aceptando un 5\,\% de falsos positivos, entrenando únicamente con tráfico benigno. Es la lectura operativa de la tabla~\ref{tab:anomalia}, que reporta el área bajo la curva.}
  \label{tab:anexo-5fpr}
  \begin{tabular}{lccc}
    \toprule
    \textbf{Régimen} & \textbf{Payload (IF)} & \textbf{Payload (AE)} & \textbf{Conducta (IF)} \\
    \midrule
    Cifrado (14-02, SSH) & 0.0000 & 0.0000 & 1.0000 \\
    En claro (22-02, Web) & 0.2808 & 0.0739 & 0.5714 \\
    Volumétrico (16-02, DoS) & 0.0000 & 0.0000 & 0.0170 \\
    \bottomrule
  \end{tabular}
\end{table}


\section{Búsqueda de hiperparámetros}

La sección~\ref{sec:sota-contraste} resume esta búsqueda en dos frases porque su
conclusión es que no cambia nada. La tabla~\ref{tab:anexo-hiperparametros}
recoge los cuatro casos completos, incluido el coste, que también informa: la
búsqueda sobre el perceptrón multicapa en el día volumétrico tarda 2.624
segundos frente a los 38 de XGBoost sobre el mismo problema, es decir, casi
setenta veces más para una ganancia de $0.0063$ frente a $0.0005$.

La búsqueda cubre los dos modelos tabulares y se valida únicamente dentro del
día de entrenamiento. El día de prueba no participa en la elección, porque
ajustar mirando el conjunto de prueba es precisamente el defecto que esta
memoria denuncia en otros sitios.

% @fuente [R.HIPERPARAMETROS[c][m][k] for c in R.HIPERPARAMETROS for m in R.HIPERPARAMETROS[c] for k in ('por_defecto', 'mejor', 'segundos')]
\begin{table}[htbp]
  \centering
  \caption[Búsqueda de hiperparámetros]{Búsqueda aleatoria de 30 configuraciones, validada únicamente dentro del día de entrenamiento. Macro-F1 con los valores por defecto, con los mejores encontrados, y coste de la búsqueda.}
  \label{tab:anexo-hiperparametros}
  \begin{tabular}{llccc}
    \toprule
    \textbf{Caso} & \textbf{Modelo} & \textbf{Por defecto} & \textbf{Mejor} & \textbf{Segundos} \\
    \midrule
    22-02 Web, metadatos & MLP & 0.9107 & 0.9383 & 37 \\
     & XGBoost & 0.9556 & 0.9679 & 46 \\
    \midrule
    15-02 DoS, metadatos & MLP & 0.9556 & 0.9619 & 2.624 \\
     & XGBoost & 0.9814 & 0.9819 & 38 \\
    \bottomrule
  \end{tabular}
\end{table}

